\input{Iteration-V2/Chapter2/chapter-02-23}

\subsection{Revision 24: key lifecycle and rollback concepts}

Key-management guidance treats cryptographic keys as lifecycle-governed objects rather than timeless constants. NIST SP 800-57 Part 1 Revision 5 discusses general key-management practices and protection across lifecycle phases \cite{nist80057r5}. Privacy Lens uses that source only as security-engineering context: neither NIST guidance nor the project's separation of issuer and approver creates a GDPR requirement or certifies a reviewer.

The Update Framework specifies version, expiry, trusted-root, and rollback-resistance concepts for software-update metadata \cite{tufspec}. Revision 24 borrows the narrower ideas of signed metadata, expiry, monotonic versions, and predecessor continuity, but does not implement TUF roles, threshold signatures, repository metadata, mirrors, consistent snapshots, or a transparency service. Calling the envelope ``TUF compliant'' would therefore exceed the evidence. The academically relevant point is the explicit translation and limitation of external security concepts, not terminological association.


\subsection{Revision 25: attestation subject and digest concepts}

The in-toto Attestation Framework separates an identified subject and digest from a predicate containing claims about that subject \cite{intotoattestation}. Revision 25 borrows the narrower discipline of binding a signed statement to an exact digest and declared release identity. It does not implement the in-toto Statement or envelope formats, a published predicate type, supply-chain layout verification, a transparency log, or in-toto conformance. The receipt-set digest is an order-independent local comparison identifier, not an inclusion proof.

This distinction matters because witness signatures can otherwise invite an unjustified inference that a binary or human role has been independently verified. Privacy Lens receipts currently identify the application by package name, semantic version, version code, and controlled-research channel, but do not include APK, AAB, or signing-certificate digests. Consequently, they can support exact-envelope agreement only; artifact identity and real-world witness governance remain separate acceptance obligations.


\subsection{Revision 26: usable-security evidence and preregistration}

Security-warning research demonstrates why a visually prominent caveat is not automatically understood. Sunshine et al. empirically compared browser warnings and showed that warning design and task context affect user response \cite{sunshine2009cryingwolf}. Bravo-Lillo et al. tested interface ``attractors'', including temporary inhibition of risky action, and found that interaction structure could redirect attention to salient information \cite{bravolillo2013attention}. Privacy Lens borrows only the hypothesis that a brief, relevant interruption may improve attention. Different domain, consequence, population, and task conditions prevent direct effect-size transfer.

The Revision-26 study draft follows the preregistration principle of recording a time-stamped, read-only study plan before collection or analysis \cite{osfpreregistration2026}. This is method preparation rather than a registered study: ethics determination, consent materials, recruitment, registration, data collection, and independent analysis remain future actions.
